% from ega4-errata-addenda-list3-fr.tex
\oldpage[IV]{359}
\begin{proof}
La nécessité de la condition est évidente par transitivité des images
réciproques de diviseurs, tous les morphismes considérés étant plats
{\hyperlink{ega4.c04.statement.21-4-4}{(21.4.4)}}. Pour tout ouvert
$U$ de $Y$, notons $f_U$ la restriction $f^{-1}(U)\to U$ de $f$; en
vertu de {\hyperlink{ega4.c04.prop.21-4-8}{(21.4.8)}} l'application
\[
 D_U\longmapsto f_U^*(D_U)
\]
de $\operatorname{Div}(U)$ dans $\operatorname{Div}(f^{-1}(U))$ est
injective; donc, si $U_1$, $U_2$ sont deux ouverts de $Y$, et si
\[
 D'|f^{-1}(U_1)=f_{U_1}^*(D_1),\qquad
 D'|f^{-1}(U_2)=f_{U_2}^*(D_2),
\]
on a nécessairement
$D_1|U_1\cap U_2=D_2|U_1\cap U_2$. Ceci montre aussitôt qu'il existe
un plus grand ouvert $U\subset Y$ tel que
\[
 D'|f^{-1}(U)=f_U^*(D_U),
 \qquad D_U\in\operatorname{Div}(U);
\]
il s'agit de voir que $U=Y$. Supposons le contraire, et soit $y$ un
point maximal de $Y-U$; on va montrer qu'il existe un voisinage ouvert
$V$ de $y$ dans $Y$ tel que
$D'|f^{-1}(V)=f_V^*(D_V)$ pour un $D_V\in\operatorname{Div}(V)$, ce qui
contredira le caractère maximal de $U$.

\textup{I) Cas où $f$ est quasi-compact.} Appliquant le fait que $f$
est quasi-compact et {\hyperlink{ega4.c04.prop.20-3-8}{(20.3.8)}}
suivant la méthode de ({\hyperref[canon-index.ega4.c04.statement.8-1-2]{8.1.2}},
a)), on est ramené à prouver qu'avec les notations de l'énoncé (mais
sans supposer d'abord de limitation sur
$\operatorname{prof}(\sh{O}_{Y,y})$), $D'_1$ est de la forme
$f_1^*(D_1)$. Or, si $\operatorname{prof}(\sh{O}_{Y,y})\leq1$, cette
assertion n'est autre que l'hypothèse. Supposons donc
$\operatorname{prof}(\sh{O}_{Y,y})\geq2$; puisque $y$ est point maximal
de $Y-U$, l'ouvert $W=Y_1\cap U$ de $Y_1$ est égal à $Y_1-\{y\}$, et
on peut se borner au cas où $Y=Y_1$, $U=Y-\{y\}$. Par hypothèse, il
existe un diviseur $D_U\in\operatorname{Div}(U)$ tel que
$f_U^*(D_U)=D'|f^{-1}(U)$; $D_U$ est classe d'équivalence d'un couple
$(\sh{L}_U,s)$, où $\sh{L}_U$ est un $\sh{O}_U$-Module inversible et
$s$ une section méromorphe régulière de $\sh{L}_U$ au-dessus de $U$
{\hyperlink{ega4.c04.prop.21-2-11}{(21.2.11)}}; par suite
{\hyperlink{ega4.c04.statement.21-4-1}{(21.4.1)}},
$D'|f^{-1}(U)$ est classe d'équivalence du couple
$(\sh{L}'_U,s')$, où
\[
 \sh{L}'_U=f_U^*(\sh{L}_U),\qquad s'=s\circ f_U.
\]
Notons $i:U\to Y$, $j:f^{-1}(U)\to X$ les injections canoniques, et
posons
\[
 \sh{L}=i_*(\sh{L}_U),\qquad
 \sh{L}'=j_*(\sh{L}'_U).
\]
Comme $f$ est plat, $f^*(\sh{L})$ est canoniquement isomorphe à
$\sh{L}'$ {\hyperref[canon-index.ega4.c04.lem.2-3-1]{(2.3.1)}}; d'autre part,
puisque $f$ est plat, la relation
$\operatorname{prof}(\sh{O}_{Y,y})\geq2$ entraîne
$\operatorname{prof}(\sh{O}_{X,x})\geq2$ pour tout $x\in f^{-1}(y)$
{\hyperref[canon-index.ega4.c04.prop.6-3-1]{(6.3.1)}}.

Or, $D'$ est classe d'équivalence d'un couple $(\sh{L}'',s'')$, où
$\sh{L}''$ est un $\sh{O}_X$-Module inversible et $s''$ une section
méromorphe régulière de $\sh{L}''$ au-dessus de $X$; $\sh{L}'_U$ est
par suite isomorphe à $\sh{L}''|f^{-1}(U)$, et il résulte alors de ce
qui précède et de {\hyperref[canon-index.ega4.c04.prop.5-9-8]{(5.9.8)}} et
{\hyperref[canon-index.ega4.c04.thm.5-10-5]{(5.10.5)}} que $\sh{L}'$ est un
$\sh{O}_X$-Module inversible isomorphe à $\sh{L}''$. Puisque $f$ est
fidèlement plat et quasi-compact, on en conclut que $\sh{L}$ est un
$\sh{O}_Y$-Module inversible
{\hyperref[canon-index.ega4.c04.prop.2-5-2]{(2.5.2)}} et par suite isomorphe à
$\sh{O}_Y$, puisque $Y$ est un schéma local. D'ailleurs, puisque
$\operatorname{prof}(\sh{O}_{Y,y})\geq1$, $U=Y-\{y\}$ est
schématiquement dense dans $Y$
({\hyperlink{ega4.c04.statement.20-2-13}{20.2.13}}, (iv)) et par suite
la section méromorphe régulière $s$ de $\sh{L}_U$ au-dessus de $U$ se
prolonge en une section méromorphe régulière $s_0$ de $\sh{L}$
au-dessus de $Y$ {\hyperlink{ega4.c04.prop.20-2-11}{(20.2.11)}}; la
classe d'équivalence de $(\sh{L},s_0)$ est donc un diviseur $D$ sur
$Y$; et il est clair que $f^*(D)$ et $D'$ induisent le même diviseur sur
$f^{-1}(U)=X-f^{-1}(y)$; mais comme
$\operatorname{prof}(\sh{O}_{X,x})\geq2$ pour tout $x\in f^{-1}(y)$,
on a $D'=f^*(D)$ {\hyperlink{ega4.c04.prop.21-1-8}{(21.1.8)}}, ce qui
achève la démonstration dans ce cas.

\textup{II) Cas où $f$ est ouvert.} Comme il s'agit de prouver
l'existence d'un voisinage ouvert $V$ de $y$ dans $Y$ tel que
$D'|f^{-1}(V)=f_V^*(D_V)$ pour un $D_V\in\operatorname{Div}(V)$, on
peut déjà supposer $Y$ affine. Puisque $f$ est ouvert, les images
$f(W)$, où $W$ parcourt les ouverts affines de $X$, forment un
recouvrement ouvert de $Y$, donc, puisque $Y$ est quasi-compact, il
existe un ouvert quasi-compact $T\subset X$ tel que $f(T)=Y$, et le
morphisme $f|T$ est quasi-compact (I, 1.1.1); \emph{a fortiori}, pour
tout ouvert quasi-compact $T_\alpha\supset T$ de $X$,

